\section{Repaso práctico dos: lista de trabajo para llevar la seguridad y la gobernanza a la práctica de ingeniería}

La mayor parte de la controversia que generó OpenClaw es, en esencia, «cómo gobernar un sistema de alta autonomía». Este capítulo convierte los principios de seguridad mencionados en la entrevista en una lista ejecutable.

\subsection{Estratificación de permisos: dividir las capacidades por nivel de riesgo}

Colocar todas las acciones en la misma capa de permisos es el error de diseño más común. Deben estratificarse por riesgo y establecer políticas de confirmación distintas.

\begin{knowledgebox}{Modelo de permisos de cuatro niveles sugerido}
\begin{itemize}
    \item \textbf{L1, operaciones de lectura}: leer páginas web públicas, bases de conocimiento, archivos no sensibles;
    \item \textbf{L2, operaciones de escritura de bajo riesgo}: borradores, etiquetas, recordatorios, archivos temporales;
    \item \textbf{L3, ejecución de riesgo medio}: enviar mensajes, abrir tickets, modificar configuraciones;
    \item \textbf{L4, ejecución de alto riesgo}: pagos, eliminaciones, publicaciones, envío externo de datos sensibles.
\end{itemize}
\end{knowledgebox}

Las de L3/L4 requieren confirmación explícita y conservación de evidencia de la operación. Si el sistema admite una «aprobación automática», también debe tener una ventana de tiempo y un mecanismo de reversión.

\subsection{Política de modelos: no es «el modelo más fuerte», sino «el modelo que mejor se ajusta»}

Que la entrevista mencione «no confíes ciegamente en los modelos baratos» no significa usar solo el modelo más caro, sino estratificar las tareas:
\begin{itemize}
    \item la recuperación y el borrador de bajo riesgo pueden usar un modelo económico;
    \item las decisiones de alto riesgo y las acciones de envío externo usan un modelo de alta confiabilidad;
    \item antes y después de las operaciones clave, hacer una verificación de consistencia entre dos modelos o una revisión por reglas.
\end{itemize}

\begin{warningbox}{El costo oculto de los modelos baratos}
Si la tasa de error de un modelo barato aumenta en escenarios de alto riesgo, el retrabajo manual posterior, los incidentes de seguridad y la pérdida de usuarios consumirán toda la ventaja de costo. El costo debe calcularse por «toda la cadena», no por el precio unitario del token.
\end{warningbox}

\subsection{Observabilidad: convertir el «comportamiento autónomo» en «comportamiento rastreable»}

\begin{importantbox}{Evidencia de ejecución que debe registrarse}
\begin{enumerate}
    \item Evidencia de entrada: mensajes, archivos, fuentes web;
    \item Evidencia de razonamiento: resumen de decisiones clave y estado de las restricciones activadas;
    \item Evidencia de ejecución: parámetros de las llamadas a herramientas, resultados devueltos, códigos de error;
    \item Evidencia de cambios: diferencias de configuración, actualizaciones de reglas, historial de reversiones.
\end{enumerate}
\end{importantbox}

Sin una cadena de evidencia, cualquier «error del agente» solo puede atribuirse a un «fallo aleatorio del modelo», lo cual es inaceptable en producción.

\subsection{Respuesta a incidentes: tratar el «accidente del modelo» como un accidente de sistema}

\begin{longtable}{p{0.2\textwidth}p{0.32\textwidth}p{0.38\textwidth}}
\toprule
\textbf{Etapa} & \textbf{Objetivo} & \textbf{Acción} \\
\midrule
Detección & Identificar el comportamiento anómalo lo antes posible & reglas de alerta, umbrales de anomalía, comparación con la línea base de comportamiento \\
\midrule
Aislamiento & Bloquear el daño continuo & degradar a modo de solo lectura, congelar las herramientas de alto riesgo \\
\midrule
Localización & Reconstruir la ruta del accidente & repetir los registros, comparar los cambios de configuración, reconstruir el contexto \\
\midrule
Corrección & Restablecer la estabilidad del servicio & revertir la política, publicar el parche, restringir temporalmente los permisos \\
\midrule
Repaso & Evitar que se repitan problemas similares & agregar casos de prueba, revisar el documento de políticas, actualizar el runbook \\
\bottomrule
\end{longtable}

\subsection{Cómo equilibrar la «ansiedad por la seguridad» y la «velocidad de construcción»}

La postura viable que ofrece la entrevista es: reconocer que la ansiedad es razonable, pero sin caer en la parálisis. La ruta más pragmática es «divulgar continuamente los riesgos, corregir continuamente las rutas y reducir continuamente el blast radius».

\begin{figure}[H]
\centering
\includegraphics[width=0.9\textwidth]{frames/frame-031210.jpg}
\caption{Segmento de cierre de la entrevista: la discusión sobre el riesgo social junto con el beneficio real para los usuarios\protect\footnotemark}
\end{figure}
\footnotetext{Intervalo de tiempo en el video: 03:12:05--03:12:20.}

\subsection{Resumen de la sección}
La seguridad no es un «punto de verificación previo al lanzamiento», sino una capacidad de operación cotidiana del producto de agente. La lección de OpenClaw es: cuanto antes se incorpore la seguridad al eje principal del producto, mayor será la capacidad de mantener el control durante la etapa de difusión.
